\section{The volume cost of the rank tests}
\label{sec:rank-root-normalization}

The replicated rank tests require the correct power of the coupling in
both the numerator and the partition function. We prove the finite-torus
estimate with that power. Write $D=\dim_{\mathbb R}\mathfrak g$, $\mu_G$
for probability Haar measure, and $\chi(u)$ for the normalized chord of
$u\in G$ used in the action. Thus $\chi(1)=0$, it is unchanged by
inversion and conjugation, and $\chi(uv)\le\chi(u)+\chi(v)$.

\subsection{Haar bounds in the physical dimension}

There are constants $0<c_-\le1\le c_+<\infty$, depending only on $G$,
such that
\begin{align}
 \mu_G\{\chi(u)\le r/4\}&\ge c_-r^D &&(0<r\le1),
       \label{eq:rank-haar-low}\\
 \int_G e^{-\chi(u)/t}\,d\mu_G(u)&\le c_+t^D &&(t>0).
       \label{eq:rank-haar-high}
\end{align}
To see this, use the actual Haar chart of Section~\ref{sec:haar-chart}.
Its density is bounded above and below by positive constants near zero.
Differentiability of the represented exponential, together with the
normalization of the colour norm, gives, after shrinking the chart,
\[
 \tfrac12\|z\|\le\chi(\exp z)\le2\|z\|.
\]
A sufficiently small chord ball is contained in that chart: the faithful
matrix embedding induces the topology of the compact group. Comparison
with Euclidean balls of radii $r/2$ and $3r$ therefore gives
$a r^D\le\mu_G\{\chi\le r\}\le b r^D$ for $0<r\le r_0$.
For $r\in[r_0,1]$ the lower bound extends by monotonicity after decreasing
$a$. The upper bound extends to all $r>0$ by the total mass bound one,
after increasing $b$.

For $k\ge0$, on $kt\le\chi(u)<(k+1)t$ the integrand in
\eqref{eq:rank-haar-high} is at most $e^{-k}$. Summing over these
measurable shells bounds the integral by
\[
 b t^D\sum_{k=0}^\infty e^{-k}(k+1)^D.
\]
The series converges, since its successive-term ratio tends to $e^{-1}$.
This proves the upper bound; rescaling the lower bound by $1/4$ proves
\eqref{eq:rank-haar-low}. The proof also covers $D=0$: powers of positive
radii are then one. No connectedness of $G$ is needed.

\subsection{An axial tree and the four unintegrated links}

Let $\Lambda=(\mathbb Z/L\mathbb Z)^4$, $L>1$, and $V=|\Lambda|$.
Represent each coordinate by $0,\ldots,L-1$. For $x\ne0$ let $p(x)$
be its largest nonzero coordinate index. Join $x$ to $x-e_{p(x)}$.
The sum of the coordinate representatives strictly decreases on every
such edge, so these edges form a tree rooted at zero. The tree has
$V-1$ positive links; there are $4V$ positive links altogether. Thus
there are $3V+1$ free links after fixing the tree links to one.

For completeness, tree gauge fixing preserves the integral of a
gauge-invariant integrable function. Integrate successively from the
leaves of the tree. At a leaf, use its tree link as the gauge variable
at that vertex and multiply all its other incident links by that
variable or its inverse. Conditional on the other links, these changes
preserve their Haar measures. The tree link becomes one, and the
remaining gauge variable integrates to one. Induction removes all
$V-1$ tree variables. This proves the identity with product probability
Haar measure on the free links, including for a disconnected group.

Keep also the four links $(-e_i,i)$, one wrap link in each direction.
They are free, distinct, and are not in the tree. We now associate a
different plaquette to each of the other $3V-3$ free links so that their
integrals can be performed successively. Here is an explicit ordering,
which also deals with wrap links. Set $p(0)=0$, and for a positive
link $e=(x,i)$ put
\[
 t(e)=\max(i,p(x)),\qquad
 \pi(e)=\bigl(t(e),\,\mathbf1_{\{i=t(e)\}},\,
                   \textstyle\sum_jx_j\bigr),
\]
ordered lexicographically.

If $i<t(e)=t$, choose the plaquette based at $x-e_t$ in directions
$(i,t)$. Its $i$-link at $x$ is $e$. Its two $t$-links are tree links,
because $x_t\ne0$ and every coordinate above $t$ vanishes. Its remaining
$i$-link has smaller tier or the same tier and smaller coordinate sum.
It therefore has smaller priority whenever it is free.

If $i=t(e)$ and $e$ is free, then $x_i=L-1$ and all coordinates above
$i$ vanish: otherwise it is a tree link. Put $w=x+e_i$. Since $e$ is
not the retained wrap link, $w\ne0$. Let $j=p(w)<i$ and choose the
plaquette based at $x-e_j$ in directions $(j,i)$. Its $i$-link at $x$
is $e$. The $j$-link based at $w-e_j$ is a tree link. The other
$j$-link has smaller tier or has the same tier with second priority
coordinate zero. The remaining $i$-link has the same tier and second
coordinate but smaller coordinate sum, since $x_j\ne0$. Thus again
every other free link on this plaquette has smaller priority.

The chosen plaquettes are distinct. If two links shared a chosen
plaquette, each would have smaller priority than the other. Moreover,
a maximal-priority remaining link occurs in no plaquette chosen for
another remaining link. With all other links fixed, its own plaquette
holonomy has the form $aub$ or $au^{-1}b$, so its distribution is Haar.
Fubini and descending induction therefore show that for continuous
$F:G\to[0,1]$,
\begin{equation}\label{eq:rank-plaquette-elimination}
 \int \prod_{p\text{ plaquette}}F(U(\partial p))\,d\mu(U)
       \le\left(\int_GF\,d\mu_G\right)^{3V-3}.
\end{equation}
Indeed discard all unchosen factors, which are at most one, and then
integrate the chosen factors in that order. The four edges of a
plaquette are distinct for $L>1$, including $L=2$; thus each integration
uses its selected variable exactly once, possibly inverted.

\subsection{The normalized estimate}

Consider a lawful representation $R$ on this torus, with gauge-invariant
path action and activities, and with nonnegative density,
coupling $\beta$, threshold $\theta\le1$, scalar exponent $\sigma$, and
clip constants satisfying $0<c_{\rm lin}<c_{\rm dir}/16$.
The path and activity bounds below are the bounds of that same
representation. These assumptions ensure that its density and the
action bounds used in tree gauge fixing are gauge invariant. Put
\[
 g=\beta^{-1/2},\quad 0<g\le\theta,\quad
 c=\min(1,c_{\rm lin}),\quad
 b(g)=N_{\mathcal A}\bigl(c_{\rm dir}g^2/2+(Ag)^2/2\bigr).
\]
Here $A$ bounds the linear curvature rows, and $N_{\mathcal A}$ is the
number of positive plaquette types. Let $r$ be the per-site bound for
the additive remainder and the cubic, quartic and loop terms. Those
bounds, with their actual support weights, give
\begin{align}
 \operatorname{exponent}_R(U)&\ge\sigma-V\{r+\beta b(g)\}
       &&\text{if every link has chord at most }g/4,
       \label{eq:rank-action-low}\\
 \operatorname{exponent}_R(U)&\le\sigma+V(r+N_{\mathcal A})
                  -\sum_p\frac{\chi(U(\partial p))}{g/c}.
       \label{eq:rank-action-high}
\end{align}
Here is the local verification of these two uses of the path bounds.
On the small-link set, every plaquette has chord at most $g$, all flags
vanish, and each linear curvature row has norm at most $Ag$ by its
absolute row bound. Summing its squared contribution and the direct
plaquette contribution gives $Vb(g)$. The other terms have absolute
sum at most $Vr$.
For the upper estimate put $a=\chi(U(\partial p))$. The clipped direct
term obeys
\[
             ca/g-1\le\beta\,\operatorname{clipped}(a).
\]
In the quadratic region this follows from
$c(a/g-1)^2\ge0$, $c\le1$, and $c\le c_{\rm dir}$.
In the linear region use $c\le c_{\rm lin}$, $g\le\theta$ and
$\beta g^2=1$. The other quadratic row term is nonnegative.
The sum over plaquettes is unchanged by signed motions; hence averaging
over them preserves this lower bound on the path action. Subtracting
the remaining action and additive terms at their absolute bound gives
\eqref{eq:rank-action-high}.

Gauge-fix the axial tree and restrict its $3V+1$ free links to
$\{\chi\le g/4\}$. The tree links are one, so
\eqref{eq:rank-action-low} applies. All flags vanish there, and the
forced gas equals one. Nonnegativity on the complement and
\eqref{eq:rank-haar-low} yield the strictly positive bound
\begin{equation}\label{eq:rank-partition}
 Z_R:=\int R\,d\mu\ge
 e^{\sigma-V(r+\beta b(g))}(c_-g^D)^{3V+1}.
\end{equation}
Now let $f$ be an integrable complex function, $z\ge0$, and suppose
\[
 z|f(U)|\le e^{\operatorname{exponent}_R(U)+VX}.
\]
Using \eqref{eq:rank-action-high}, then
\eqref{eq:rank-plaquette-elimination} and
\eqref{eq:rank-haar-high}, gives
\[
 z\left|\int f\,d\mu\right|
 \le e^{\sigma+V(r+N_{\mathcal A}+X)}
             \{c_+(g/c)^D\}^{3V-3}.
\]
Define $h=\log c_+-D\log c\ge0$ and
\[
 P_R(X)=2r+N_{\mathcal A}+\beta b(g)+X+3h-4\log c_-.
\]
Division by \eqref{eq:rank-partition} gives
\begin{equation}\label{eq:rank-ratio}
 \frac{z|\int f\,d\mu|}{Z_R}
       \le \exp\{VP_R(X)-4D\log g\}.
\end{equation}
To check the constants, the powers of $g^D$ differ by exactly four.
Also $(3V-3)h\le3Vh$ and
$-(3V+1)\log c_-\le-4V\log c_-$, since $V\ge1$ and
$\log c_-\le0$. These are the only enlargements in this division.

The same estimate holds for the auxiliary history of
Section~\ref{sec:restarted-rank-history}. If its additive terms have
value bound $\delta$ per site, delete them temporarily to obtain $R^0$.
Then $e^{-V\delta}R^0\le R\le e^{V\delta}R^0$ because $R\ge0$,
and $R^0=e^ER\ge0$. Apply the preceding estimate to $R^0$ with
numerator expense $X+\delta$, and use
$Z_{R^0}\le e^{V\delta}Z_R$. The extra numerator and denominator
expenses total $2V\delta$. This is exactly the contribution of
$\delta$ to $2r$ in $P_R$. Derivative estimates on those auxiliary
additive terms are not needed.

\subsection{Taking the replica root}

Let a terminal tiling have $q$ cells, and push its full marked density
to a torus with exactly $q$ sites along the auxiliary history. The
unnormalized scalar integral remains the original partition function
$Z$. Suppose each coefficient of the pushed marked density has the
local estimate needed above with coefficient price $a^{q|I|}$ and
per-site expense $X\le1+M$, where $M$ is the maximum number of markers
per cell and $I\ne\varnothing$ is the chosen marker subset. The signed
orientation copies use conjugation on the reversed cells. Thus the
coefficient carrying $I$ in every cell is exactly $Z$ times the
replicated Wilson expectation. Equation~\eqref{eq:rank-ratio} gives
\[
 a^{q|I|}|\mathbb E(\text{replicated coefficient})|
       \le e^{qP_R(X)-4D\log g}.
\]
Taking the positive $q$th root proves
\begin{equation}\label{eq:rank-root}
 \|\text{coefficient }I\|_{\rm root}
 \le\frac{\exp\{P_R(X)-4D\log g/q\}}{a^{|I|}}.
\end{equation}
This retains the finite-volume harmonic correction $-4D\log g/q$.

The clock buys this last expense uniformly in the period. At an output
level $m$, $\beta\le4\widehat\beta_m$, so
$-4D\log g=2D\log\beta\le8D\widehat\beta_m$.
The clock coupling grows linearly in its logarithmic spatial scale,
whereas the minimal torus site count grows at least exponentially in
that scale. Consequently $8D\widehat\beta_m/q\le1$ eventually, even
at the smallest period. For the remaining terms, use the same clip and row bound $A$ as the
clock stage, and suppose
\[
 \theta_R\le\theta_{\rm clock},\qquad
 \Lambda\le\Lambda_{\rm clock},\qquad \delta\le1.
\]
The late-scale purchases are
\[
 \widehat\beta_m\theta_{\rm clock}^{3}\le1,\qquad
 \widehat\beta_m\theta_{\rm clock}^{4}\le1,\qquad
 \Lambda_{\rm clock}\theta_{\rm clock}^{2}\le1.
\]
Consequently the actual remainder bound satisfies
\[
 r=\delta+\beta(A\theta_R^3+A\theta_R^4)
                     +\Lambda\theta_R^2
       \le1+8A+1,
 \qquad
 \beta b(g)=\frac{N_{\mathcal A}}2(c_{\rm dir}+A^2).
\]
With $X\le1+M$, a fixed stage ticket bounding $P_R(X)+1$ is therefore
\[
 T_M=2(1+8A+1)+N_{\mathcal A}
       +\frac{N_{\mathcal A}}2(c_{\rm dir}+A^2)
       +(1+M)+3h-4\log c_-+1.
\]
The final $1$ pays for the harmonic correction, rather than enlarging
an already spent ticket. Choose the comparison depth $k$ on a
constant-profile interval and purchase
\[
 T_M\le I_k=\frac{\log \mathcal M_k}{16M},
\]
where $\mathcal M_k$ is the diverging clock monitor. This is a
post-block-length scale delay. Set $a=e^{2I_k}$. Since $|I|\ge1$,
\eqref{eq:rank-root} is at most $e^{I_k}/a=e^{-I_k}\to0$.
For the rank detectors $M=2$. These estimates apply to the one-point
and corrected two-point coefficients themselves. The local marked
construction must supply their displayed coefficient price on the
same auxiliary density; no estimate for a different denominator is
used here.
